% Preâmbulo do caderno de estudo (LuaLaTeX). Incluído pelo pandoc via --include-in-header.
% Paleta Flexoki, a mesma de R/estilo_grafico.R.

% ---- fontes: Libertinus + reserva para símbolos que ela não tem (●, ✔, ⚠...) ----
\directlua{luaotfload.add_fallback("simbolos", {"Segoe UI Symbol:mode=harf;", "Cambria Math:mode=harf;", "Segoe UI Emoji:mode=harf;"})}
\setmainfont{Libertinus Serif}[RawFeature={fallback=simbolos}]
\setsansfont{Libertinus Sans}[RawFeature={fallback=simbolos}]
\setmonofont{Consolas}[Scale=0.84, RawFeature={fallback=simbolos}]
\setmathfont{Libertinus Math}

% macros do MathJax que o LaTeX não conhece
\providecommand{\gt}{>}
\providecommand{\lt}{<}

% ---- cores ----
\usepackage{xcolor}
\definecolor{fxpapel}{HTML}{FFFCF0}
\definecolor{fxtexto}{HTML}{100F0F}
\definecolor{fxeixo}{HTML}{6F6E69}
\definecolor{fxgrade}{HTML}{E6E4D9}
\definecolor{fxazul}{HTML}{205EA6}
\definecolor{fxverde}{HTML}{66800B}
\definecolor{fxteal}{HTML}{24837B}
\definecolor{fxroxo}{HTML}{5E409D}
\definecolor{fxlaranja}{HTML}{BC5215}
\definecolor{fxvermelho}{HTML}{AF3029}
\definecolor{fxamarelo}{HTML}{D0A215}

% ---- callouts do Obsidian/GitHub ----
\usepackage[most]{tcolorbox}
\newtcolorbox{callout}[2]{%
  enhanced, breakable, colback=#1!5!white, colframe=#1, boxrule=0pt,
  leftrule=2.6pt, arc=1.5pt, outer arc=1.5pt,
  left=7pt, right=7pt, top=5pt, bottom=5pt, before skip=8pt, after skip=8pt,
  fonttitle=\sffamily\bfseries\small, coltitle=#1!80!black, colbacktitle=#1!12!white,
  titlerule=0pt, toptitle=2pt, bottomtitle=2pt,
  title={#2}}

% ---- código: quebra de linha longa e fundo discreto ----
\usepackage{fvextra}
\fvset{breaklines=true, breakanywhere=true, fontsize=\small}
\makeatletter
\@ifundefined{Shaded}{}{\renewenvironment{Shaded}{\begin{tcolorbox}[enhanced, breakable, colback=fxgrade!35!white, boxrule=0pt, arc=1pt, left=4pt, right=4pt, top=3pt, bottom=3pt]}{\end{tcolorbox}}}
\makeatother
\RecustomVerbatimEnvironment{verbatim}{Verbatim}{breaklines=true, breakanywhere=true, fontsize=\small}

% ---- títulos ----
\usepackage{titlesec}
\ifdefined\chapter
\titleformat{\chapter}[display]
  {\normalfont\sffamily\bfseries\color{fxtexto}}
  {\small\mdseries\color{fxteal}\MakeUppercase{\chaptertitlename}~\thechapter}
  {3pt}{\LARGE\raggedright}
\titlespacing*{\chapter}{0pt}{-12pt}{16pt}
\fi
\titleformat{\section}{\normalfont\sffamily\bfseries\Large\color{fxtexto}\raggedright}{\thesection}{0.6em}{}
\titleformat{\subsection}{\normalfont\sffamily\bfseries\large\color{fxtexto}\raggedright}{\thesubsection}{0.6em}{}
\titleformat{\subsubsection}{\normalfont\sffamily\bfseries\normalsize\color{fxeixo}\raggedright}{\thesubsubsection}{0.6em}{}

% ---- cabeçalho e rodapé ----
\usepackage{fancyhdr}
\setlength{\headheight}{22pt}
\pagestyle{fancy}
\fancyhf{}
\fancyhead[L]{\sffamily\footnotesize\color{fxeixo}\nouppercase{\leftmark}}
\fancyhead[R]{\sffamily\footnotesize\color{fxeixo}Econometria I · P1}
\fancyfoot[C]{\sffamily\footnotesize\color{fxeixo}\thepage}
\renewcommand{\headrulewidth}{0.3pt}
\renewcommand{\headrule}{\hbox to\headwidth{\color{fxgrade}\leaders\hrule height \headrulewidth\hfill}}
\fancypagestyle{plain}{\fancyhf{}\fancyfoot[C]{\sffamily\footnotesize\color{fxeixo}\thepage}\renewcommand{\headrulewidth}{0pt}}
\ifdefined\chapter
\renewcommand{\chaptermark}[1]{\markboth{#1}{#1}}
\else
\fancyhead[L]{\sffamily\footnotesize\color{fxeixo}Simulado da P1}
\fi

% ---- tabelas e texto ----
\renewcommand{\arraystretch}{1.18}
\setlength{\emergencystretch}{3em}
\setlength{\parskip}{4pt plus 1pt}
\setlength{\parindent}{0pt}
% hyperref é carregado pelo template do pandoc; as cores vêm de -V linkcolor/urlcolor
